\uuid{6K6K}
\exo7id{3505}
\titre{exo7 3505}
\auteur{quercia}
\organisation{exo7}
\datecreate{2010-03-10}
\isIndication{false}
\isCorrection{true}
\chapitre{Réduction d'endomorphisme, polynôme annulateur}
\sousChapitre{Diagonalisation}
\cpge{mpsi-mp2i.changements-bases.a, pcsi.matrices-applications-lineaires.c, ptsi.matrices-applications-lineaires.c, mpsi-mp2i.changements-bases.c, bcpst1.matrices, mp-mpi.reduction-endomorphismes.b, pc.reduction.a, psi.reduction-endomorphismes.a, pt.reduction-endomorphismes.a, mp-mpi.reduction-endomorphismes.c, pc.reduction.b, psi.reduction-endomorphismes.b, pt.reduction-endomorphismes.b, mp-mpi.reduction-endomorphismes.d, pc.reduction.c, psi.reduction-endomorphismes.c, pt.reduction-endomorphismes.c, bcpst2.valeurs-propres-vecteurs-propres.a, bcpst2.valeurs-propres-vecteurs-propres.b, bcpst2.produit-scalaire.c}
\module{Algèbre}
\niveau{L2}
\difficulte{}

\contenu{
\question{
Diagonaliser la matrice
$M = \begin{pmatrix} e &a &b &c \cr
               a &e &c &b \cr
               b &c &e &a \cr
               c &b &a &e \cr\end{pmatrix} \in \mathcal{M}_4(\R)$.
}
\reponse{
$P = \begin{pmatrix}1& 1& 1& 1 \cr
                        1& 1&-1&-1 \cr
                        1&-1& 1&-1 \cr
                        1&-1&-1& 1 \cr\end{pmatrix}$,
         $D = \text{diag}\,(a+b+c+e,a-b-c+e,-a+b-c+e,-a-b+c+e)$.
}
}
